\documentclass[11pt,a4paper]{article}
\usepackage[utf8]{inputenc}
\usepackage{amsmath,amssymb,amsthm,mathtools}
\usepackage[margin=1.1in]{geometry}
\usepackage{hyperref}
\usepackage{booktabs}

\theoremstyle{plain}
\newtheorem{theorem}{Theorem}[section]
\newtheorem{lemma}[theorem]{Lemma}
\newtheorem{proposition}[theorem]{Proposition}
\newtheorem{corollary}[theorem]{Corollary}
\newtheorem{conjecture}[theorem]{Conjecture}
\theoremstyle{definition}
\newtheorem{definition}[theorem]{Definition}
\newtheorem{remark}[theorem]{Remark}

\newcommand{\R}{\mathbb{R}}
\newcommand{\diag}{\operatorname{diag}}
\newcommand{\tr}{\operatorname{tr}}

\title{Congruence Rigidity and the Fusion Bound:\\
what a positive weight can and cannot do\\
to the spectrum of a transfer kernel}
\author{Lluis Eriksson}
\date{Draft --- 1 August 2026}

\begin{document}
\maketitle

\begin{abstract}
A positive site weight acts on a transfer kernel by \emph{congruence},
$K \mapsto DKD$ with $D$ positive diagonal, not by similarity. Similarity
preserves the entire spectrum; congruence preserves strictly less, and strictly
more than nothing. We determine both halves for the kernel of $L$ decoupled
Ising bonds. \emph{Rigid half:} the sign of every quadratic form value survives,
so definiteness in either sign is a congruence invariant --- the definite case
of Sylvester's law of inertia, recorded here in the elementary form that a
machine can check without an eigenvalue count. \emph{Fragile half:} the value of
the spectral ratio $r = |\lambda_1|/|\lambda_0|$ does not survive at all, and we
locate exactly how far it moves. Restricting the $L$-site kernel to the two
antipodal configurations leaves \emph{one} Ising bond of coupling $\beta L$; its
ratio is $\tanh(\beta L)$. A weight concentrating on the ferromagnetic
configurations therefore \emph{fuses} the $L$ sites into a single effective site
carrying $L$ times the coupling. Since $\tanh(\beta L)\to 1$, no bound on $r$ can
be uniform in $L$: the obstruction to volume-uniformity is named, exactly, and
it is a property of the congruence and not of any spectral estimate.
Consequently an argument that bounds $r$ using only congruence invariants ---
rank, inertia, definiteness --- cannot produce an $L$-uniform bound at all.
We then ask which congruence is extremal. For \emph{uniform} concentrations the
answer is proved: keeping $m$ sites of mutual correlation $\mu$ leaves
$\mu J+(1-\mu)I$, of ratio $(1-\mu)/(1+(m-1)\mu)$, which is strictly decreasing
in $m$ --- the more sites the weight keeps, the smaller the ratio it can reach,
so the best uniform concentration is a \emph{pair}. For arbitrary weights we
record, with certified numerics and \emph{without proof}, the conjecture that
$\sup_{D>0}r(DMD)=(1-\mu)/(1+\mu)$ with $\mu=\min_{i\neq j}M_{ij}$: the extremal
congruence concentrates on the \emph{least correlated pair}. The hypercube
instance is the fusion bound. All theorems are machine-checked in Lean~4.
\end{abstract}

\section{The question}

Let $M$ be a real symmetric matrix and $D$ an invertible diagonal matrix with
positive entries. The map
\[
  M \;\longmapsto\; DMD
\]
is a \emph{congruence}. It arises whenever a statistical-mechanical transfer
kernel is symmetrised against a site weight: if the unnormalised kernel is
$w(\sigma)K(\sigma,\tau)$, the symmetric form with the same spectrum is
$\sqrt{w(\sigma)}\,K(\sigma,\tau)\sqrt{w(\tau)}$, that is $DKD$ with
$D=\diag(\sqrt{w})$.

Congruence is not similarity. Similarity $M \mapsto S M S^{-1}$ preserves the
characteristic polynomial and therefore everything spectral. Congruence
preserves the quadratic form only up to a change of variable, and the change of
variable is not an isometry. The natural question --- and, we will argue, the
one that governs whether a spectral estimate can be uniform in the system size
--- is:

\begin{center}
\emph{Which spectral data survive a positive-diagonal congruence, and how far
can the rest move?}
\end{center}

We answer both parts for the kernel of $L$ decoupled Ising bonds. Section~2
gives the rigid half, Section~3 the fragile half, Section~4 the consequence,
and Section~5 states a conjectural general form with its numerical certificate.

\paragraph{What this paper does not do.} It does not decide whether any
particular coupled transfer operator has a spectral gap uniform in its
extension. It \emph{explains} why such a question is hard by exhibiting the
exact mechanism that destroys uniformity, and explaining is not resolving. No
statement here has a consequence for gauge theory in any dimension, and none is
claimed.

\section{The rigid half: congruence cannot move a sign}

Write $q_M(x)=\sum_{i,j} x_i M_{ij} x_j$ for the quadratic form of $M$.

\begin{lemma}[Change of variable]\label{lem:change}
For every $M$, every $d$, and every $x$,
\[
  q_{\,\diag(d)\,M\,\diag(d)}(x) \;=\; q_M(dx),
  \qquad (dx)_i := d_i x_i .
\]
\end{lemma}

\begin{proof}
$(\diag(d)M\diag(d))_{ij} = d_i M_{ij} d_j$; substitute and reassociate.
\end{proof}

\begin{theorem}[Definiteness is a congruence invariant]\label{thm:rigid}
Let $d_i \neq 0$ for every $i$. Then $\diag(d)M\diag(d)$ is positive definite
if and only if $M$ is, and negative definite if and only if $M$ is.
\end{theorem}

\begin{proof}
By Lemma~\ref{lem:change} the two forms have the same value set, because
$x \mapsto dx$ is a bijection of $\R^n$ (its inverse is $x \mapsto d^{-1}x$) and
carries nonzero vectors to nonzero vectors: if $d_ix_i=0$ with $d_i\neq0$ then
$x_i=0$. Apply this in both directions.
\end{proof}

Theorem~\ref{thm:rigid} is the definite case of Sylvester's law of inertia
\cite{sylvester1852}, which says more: the full signature $(n_+,n_0,n_-)$ is a
congruence invariant. We do not reprove the general law. We isolate the definite
case because it admits a proof through the form alone --- no diagonalisation, no
eigenvalue count --- and that is what makes it cheap to certify mechanically.

\begin{remark}
Invertibility of $D$ is not decorative. A diagonal with $k$ zero entries drops
the rank by $k$ and can turn a definite form into a degenerate one. Every claim
below assumes $D$ strictly positive.
\end{remark}

\section{The fragile half: the fusion identity}

Let $s(\sigma_j,\tau_j)=+1$ if $\sigma_j=\tau_j$ and $-1$ otherwise, and let
\[
  K_\beta^{(L)}(\sigma,\tau)
    \;=\; \exp\Big(\beta \sum_{j=1}^{L} s(\sigma_j,\tau_j)\Big),
  \qquad \sigma,\tau \in \{0,1\}^L ,
\]
the kernel of $L$ decoupled Ising bonds at coupling $\beta$. Write
$\sigma^+ = (0,\dots,0)$ and $\sigma^-=(1,\dots,1)$ for the two antipodal
(ferromagnetic) configurations, and let
\[
  B(a) \;=\; \begin{pmatrix} e^{a} & e^{-a} \\ e^{-a} & e^{a}\end{pmatrix}
\]
be the transfer matrix of a \emph{single} Ising bond of coupling $a$.

\begin{lemma}[One bond]\label{lem:bond}
$B(a)$ has eigenvectors $(1,1)$ and $(1,-1)$ with eigenvalues $2\cosh a$ and
$2\sinh a$. Its spectral ratio is
\[
  \frac{2\sinh a}{2\cosh a} \;=\; \tanh a .
\]
\end{lemma}

\begin{proof}
Direct: $B(a)(1,1)^{\!\top} = (e^a+e^{-a})(1,1)^{\!\top}$ and
$B(a)(1,-1)^{\!\top} = (e^a-e^{-a})(1,-1)^{\!\top}$. Both eigenvalues are real
and $2\cosh a>0$, so the ratio is well defined.
\end{proof}

\begin{theorem}[Fusion identity]\label{thm:fusion}
The $2\times 2$ principal submatrix of $K_\beta^{(L)}$ on the antipodal pair
$\{\sigma^+,\sigma^-\}$ is exactly one Ising bond of coupling $\beta L$:
\[
  \begin{pmatrix}
    K_\beta^{(L)}(\sigma^+,\sigma^+) & K_\beta^{(L)}(\sigma^+,\sigma^-)\\[2pt]
    K_\beta^{(L)}(\sigma^-,\sigma^+) & K_\beta^{(L)}(\sigma^-,\sigma^-)
  \end{pmatrix}
  \;=\; B(\beta L),
\]
and its spectral ratio is $\tanh(\beta L)$.
\end{theorem}

\begin{proof}
$\sigma^+$ agrees with itself at all $L$ sites, so the sum of signs is $+L$ and
the entry is $e^{\beta L}$; likewise for $\sigma^-$. The two antipodal
configurations disagree at all $L$ sites, so the sum is $-L$ and the entry is
$e^{-\beta L}$. Compare with $B(\beta L)$ and apply Lemma~\ref{lem:bond}.
\end{proof}

Theorem~\ref{thm:fusion} is the mechanism this paper is about, and it is worth
stating in words. A weight that concentrates on a set $S$ of configurations does
not \emph{perturb} the kernel --- it \emph{restricts} it to $S$. When $S$ is the
antipodal pair, the restriction is a single bond whose coupling is $L$ times the
original. \textbf{The weight fuses the $L$ sites into one effective site.}

\section{Why that is a wall}

\begin{theorem}[No uniform ratio bound]\label{thm:wall}
For every $\beta>0$ and every $\rho<1$ there is an $L$ with
$\tanh(\beta L)>\rho$. Consequently the fused ratio is not bounded away from
$1$, and no estimate of the form $r \le \rho <1$ can hold uniformly in the
number of sites.
\end{theorem}

\begin{proof}
$\tanh a = 1 - 2/(e^{2a}+1)$ is strictly increasing with limit $1$; choose $L$
with $e^{2\beta L}+1 > 2/(1-\rho)$, which exists because $\beta>0$ and
$e^{x}>x$.
\end{proof}

This is the honest content of the paper's title. It is \emph{not} a proof that
any given coupled transfer operator fails to have a uniform gap; the
concentration that realises the fusion is a limit, and the weight of a physical
model may or may not approach it. What Theorem~\ref{thm:wall} establishes is
that \emph{the obstruction lives in the congruence}, not in whatever spectral
estimate one happens to be attempting. An argument that bounds $r$ using only
data preserved by congruence --- rank, inertia, definiteness --- cannot
possibly produce an $L$-uniform bound, because those data are identical for the
unfused kernel (ratio $\tanh\beta$) and for every fused one (ratio up to
$\tanh(\beta L)$).

\begin{corollary}
Any $L$-uniform ratio bound must use a quantity that a positive-diagonal
congruence can change.
\end{corollary}

\subsection*{The orbit, not its closure}

The fusion of Theorem~\ref{thm:fusion} is realised by a weight that
\emph{concentrates}, and the congruences considered here are by \emph{strictly
positive invertible} diagonals. Setting the remaining diagonal entries to zero
would leave the orbit and land in its closure, so the interface deserves to be
stated rather than assumed. It is stated, and the statement is stronger than a
limiting argument.

\begin{proposition}[Concentration stays inside the orbit]\label{prop:orbit}
Fix two configurations $p\neq q$ and $0<\varepsilon\le 1$, and let
$D_\varepsilon$ be the diagonal that is $1$ at $p$ and $q$ and $\varepsilon$
elsewhere. Then $D_\varepsilon$ is strictly positive, and for \emph{every} such
$\varepsilon$,
\[
  (D_\varepsilon M D_\varepsilon)_{pq}=M_{pq},
  \qquad
  (D_\varepsilon M D_\varepsilon)_{qp}=M_{qp},
\]
while $|(D_\varepsilon M D_\varepsilon)_{ij}|\le\varepsilon\,|M_{ij}|$ whenever
$i\notin\{p,q\}$.
\end{proposition}

\begin{proof}
$(D_\varepsilon M D_\varepsilon)_{ij}=d_id_jM_{ij}$ with $d_p=d_q=1$, giving the
two equalities; and $d_i=\varepsilon$, $d_j\le1$ off the pair.
\end{proof}

So the fused bond is \emph{already present at every strictly positive
$\varepsilon$}: the limit does not create the block, it deletes the rest of the
matrix. Converting that deletion into a statement about eigenvalues is the
classical continuity of the spectrum in finite dimensions (eigenvalues depend
continuously on the entries, e.g.\ \cite[Appendix D]{hornjohnson}); we cite that
step rather than reprove it, and Proposition~\ref{prop:orbit} is where the
citation attaches.

\section{Which congruence is extremal: pairs beat exchangeable clusters}

Theorem~\ref{thm:fusion} exhibits \emph{one} congruence reaching
$\tanh(\beta L)$. Is it the best one? For the family of \emph{uniform}
concentrations the question has a complete and elementary answer, and the answer
explains the mechanism.

\begin{theorem}[Exchangeable concentration]\label{thm:exch}
Fix $0<\mu<1$ and let $E_m=\mu J_m+(1-\mu)I_m$ be the exchangeable correlation
matrix on $m$ sites. Its spectrum consists of exactly two points: $1+(m-1)\mu$ on
the constant vector, and $1-\mu$ on every vector summing to zero. Hence the ratio
it reaches is
\[
  \frac{1-\mu}{1+(m-1)\mu},
\]
which is \emph{strictly decreasing} in $m$.
\end{theorem}

\begin{proof}
$E_m\mathbf 1 = (1+(m-1)\mu)\mathbf 1$ by summing a row. If $\sum_i x_i=0$ then
$(E_mx)_i=\mu\sum_j x_j+(1-\mu)x_i=(1-\mu)x_i$. These two eigenspaces have
dimensions $1$ and $m-1$, so they exhaust the spectrum. For monotonicity,
$k\mapsto (1-\mu)/(1+k\mu)$ is strictly decreasing on $k\ge 0$ because the
numerator is positive and the denominator strictly increasing.
\end{proof}

The reading is the point. \textbf{The more sites an exchangeable concentration
keeps, the smaller the ratio it can reach.} So among these the best is the
smallest nondegenerate one --- a \emph{pair} --- and at $m=2$ the value is
$(1-\mu)/(1+\mu)$. On the hypercube kernel, $\mu=e^{-2\beta L}$ at the antipodal
pair, and $(1-\mu)/(1+\mu)=\tanh(\beta L)$: Theorem~\ref{thm:fusion} is the
$m=2$ case of Theorem~\ref{thm:exch}.

\begin{remark}[Exactly how far this goes]\label{rem:exchscope}
Theorem~\ref{thm:exch} concerns subsets on which \emph{all} pairwise
correlations equal a common $\mu$. On a general subset the entries $M_{ij}$
differ and the spectrum has no such closed form, so the theorem does \emph{not}
say that a pair beats every uniformly weighted subset of an arbitrary $M$. The
section title should be read with the word \emph{exchangeable} attached. What
the theorem does establish is the mechanism --- spreading a concentration over
more equally correlated sites lowers the reachable ratio --- and it is that
mechanism, not a general proof, that makes Conjecture~\ref{conj:lcp} plausible.
\end{remark}

That settles the exchangeable family. For arbitrary weights we have only a
conjecture, and we say so wherever it appears.

\begin{conjecture}[Least-correlated pair]\label{conj:lcp}
Let $M$ be symmetric with $M_{ii}=1$ and $0<M_{ij}<1$ for $i\neq j$, and let
$\mu = \min_{i\neq j}M_{ij}$. Then
\[
  \sup_{D>0 \text{ diagonal}} r(DMD) \;=\; \frac{1-\mu}{1+\mu},
\]
approached by concentrating $D$ on an argmin pair, and not attained.
\end{conjecture}

The lower bound is elementary: concentrating on a pair $\{i,j\}$ leaves
$\begin{psmallmatrix}1 & M_{ij}\\ M_{ij} & 1\end{psmallmatrix}$ up to scale,
of ratio $(1-M_{ij})/(1+M_{ij})$, which is largest at the least correlated
pair. Theorem~\ref{thm:exch} adds that no \emph{uniform} concentration on more
sites can beat it. \textbf{The upper bound over arbitrary weights is open.} We
have no proof, and we say so wherever the statement appears.

For the record, the reduction we did obtain and could not close: since
$T$ is positive definite exactly when $M$ is (Theorem~\ref{thm:rigid}),
\[
  r(T)\le\frac{1-\mu}{1+\mu}
  \iff
  \lambda_0-\lambda_1\;\ge\;\mu\,(\lambda_0+\lambda_1),
\]
and writing $M=\mu J+(1-\mu)N$ --- where $N$ again has unit diagonal and
entries in $[0,1)$ --- gives $T=\mu\,dd^{\!\top}+(1-\mu)\,DND$, a rank-one
positive semidefinite piece plus an entrywise nonnegative one. Weyl's inequality
and the Perron root of the second piece bound $|\lambda_1|$ correctly but lose
too much on $\lambda_0$: testing against the Perron vector of $DND$ reduces the
claim to $\sum_{i\neq j}(1+\mu-2M_{ij})u_iu_j\ge(1-\mu)\sum_i u_i^2$, whose
left-hand side is negative once some $M_{ij}>(1+\mu)/2$. The top eigenvector of
$T$ is not the Perron vector of $DND$, and the slack is exactly that gap.

On the hypercube, $K_\beta^{(L)}$ normalised to unit diagonal has
$M_{\sigma\tau}=e^{-2\beta h(\sigma,\tau)}$ with $h$ the Hamming distance; the
least correlated pair is the antipodal one, $\mu=e^{-2\beta L}$, and
$(1-\mu)/(1+\mu)=\tanh(\beta L)$. Conjecture~\ref{conj:lcp} therefore contains
Theorem~\ref{thm:fusion} as its hypercube instance.

\paragraph{Numerical certificate.} Random search plus Nelder--Mead ascent over
positive diagonals, on families chosen to break the conjecture, never exceeded
$(1-\mu)/(1+\mu)$ and in almost every case attained it to machine precision.

\begin{center}
\begin{tabular}{lrr}
\toprule
family & worst excess over target & cases \\
\midrule
hypercube, $L\in\{2,3,4\}$, $\beta\in\{0.15,0.4,0.8\}$ & $1.0\times10^{-15}$ & 9\\
random positive Gram, $n\in\{4,5,6\}$                  & $5.0\times10^{-16}$ & 12\\
full rank / near singular ($\kappa\sim6\times10^{6}$)  & $3.3\times10^{-16}$ & 4\\
near-exchangeable, isolated argmin, entries near $1$   & $4.4\times10^{-16}$ & 3\\
exactly exchangeable $\mu J+(1-\mu)I$                  & $1.0\times10^{-15}$ & 2\\
\emph{indefinite}, $\lambda_{\min}\in[-0.42,-0.06]$    & $1.6\times10^{-15}$ & 12\\
larger sizes, $n\in\{8,10,12\}$, $\kappa$ up to $5\times10^{4}$
                                                       & $4.2\times10^{-16}$ & 6\\
\bottomrule
\end{tabular}
\end{center}

The last row was added after an audit of our own procedure: the adversarial
sweep as first written never actually reached $n=8$. Its positive-definiteness
\emph{precondition} fired on that family, the generator was repaired, and the
family was not re-run --- so the largest size with evidence had been $n=6$, not
$n=12$ as the sweep's summary implied. A precondition that is never met carries
no evidence either way, and the script now returns a distinct
\textsc{inconclusive} code rather than reporting such a case as a result.

The last row deserves comment, because it corrected us. A pre-registered gate
predicted that positive definiteness would be load-bearing --- that at least one
of twelve indefinite witnesses would break the bound. None did; all twelve
landed \emph{on} it. The hypothesis was therefore dropped, and
Conjecture~\ref{conj:lcp} as stated above assumes no definiteness. The gate, its
prediction, and its failure are recorded in the repository rather than repaired
in place.

\paragraph{Scope of the certificate.} These are numerics, not proof. They cover
$n\le 12$, $\mu>0$, and $\lambda_{\min}\ge -0.43$. Nothing here tests $\mu\le 0$
(where the target exceeds $1$ and the question changes character), nor a
negative eigenvalue large enough to make $|\lambda_0|$ the negative extreme, at
which point $r$ measures something else.

\section{Relation to prior work}

\paragraph{Sylvester (1852).} The law of inertia is classical and
Theorem~\ref{thm:rigid} is its definite case. We claim no novelty on the rigid
half; the contribution is that this form of it is checkable without
diagonalisation.

\paragraph{Ostrowski.} Ostrowski's quantitative form of the inertia law gives,
for Hermitian $A$ and invertible $S$, $\lambda_k(SAS^{*})=\theta_k\lambda_k(A)$
with $\theta_k\in[\lambda_{\min}(SS^{*}),\lambda_{\max}(SS^{*})]$. For $S=D$
diagonal this yields
\[
  r(DAD) \;\le\; \kappa(D^2)\, r(A),
  \qquad \kappa(D^2)=\frac{\max_i d_i^2}{\min_i d_i^2},
\]
which \textbf{diverges} exactly in the regime of interest, namely when $D$
degenerates. Theorem~\ref{thm:fusion} and Conjecture~\ref{conj:lcp} supply a
\emph{finite} substitute in that regime: the supremum is $\tanh(\beta L)$,
respectively $(1-\mu)/(1+\mu)$, however badly conditioned $D$ becomes. That
contrast --- divergent classical bound, finite exact value --- is the delta of
Section~5, and it is stated here side by side so that it can be checked.

\paragraph{Onsager and Kaufman.} The kernel $DK_\beta D$ with a ferromagnetic
nearest-neighbour weight is the transfer matrix of the two-dimensional
anisotropic Ising model on a strip, whose spectrum is known in closed form.
Nothing in Sections~3--4 is new mathematics about that model, and the reader who
knows the exact solution will find $\tanh(\beta L)$ unsurprising: it is the
zero-temperature limit in the transverse direction, where the strip collapses to
a single spin. The contribution is not the value. It is that the value is the
supremum \emph{within the family of uniform concentrations}
(Theorem~\ref{thm:exch}) and a lower bound for the full diagonal orbit, that the
orbit's rigid invariants are simultaneously pinned, and that the pair of facts
identifies which quantities an argument may and may not use. We are careful with
that phrasing: whether $\tanh(\beta L)$ is the supremum over \emph{all} positive
diagonals is Conjecture~\ref{conj:lcp} and is not proved here.

\paragraph{Formalisation.} We are not aware of a mechanised treatment of
congruence orbits of transfer kernels. This is an absence of found prior art in
the searches we ran, not a claim of priority.

\section{What is open}

\begin{enumerate}
\item The upper bound in Conjecture~\ref{conj:lcp} \emph{over arbitrary
      weights}. Theorem~\ref{thm:exch} settles it for uniform concentrations;
      the gap is everything else. A proof would make the fusion bound an exact
      classification rather than an attained lower bound with a certificate,
      and \S5 records where our attempt loses.
\item Whether the physical weight $w_\gamma$ attains the antipodal restriction
      in the limit $\gamma\to\infty$ with the rate we measure. The deficit
      $\tanh(\beta L)-r(\gamma)$ was found to decay like $e^{-2\gamma}$ (fitted
      exponents $-1.96$ to $-2.08$ against a pre-registered prediction of $-2$),
      consistent with one-domain-wall configurations dominating the correction.
      The limit itself is not proved here.
\item The general signature-counting law as a mechanised statement.
\end{enumerate}

\section*{Machine verification --- status}

Every theorem, lemma and proposition stated above --- the rigid half
(Theorem~\ref{thm:rigid}, Lemma~\ref{lem:change}), the one-bond spectrum
(Lemma~\ref{lem:bond}), the fusion identity (Theorem~\ref{thm:fusion}), the wall
(Theorem~\ref{thm:wall}), the orbit statement (Proposition~\ref{prop:orbit}) and
the exchangeable computation (Theorem~\ref{thm:exch}) --- is written in Lean~4
against a pinned Mathlib in \texttt{YangMills/OS/CongruenceSpectrum.lean}, and
\textbf{has been verified}. Two things are not, and both are labelled: the
continuity of eigenvalues invoked after Proposition~\ref{prop:orbit}, which is
classical and cited; and Conjecture~\ref{conj:lcp}, which carries the word
\emph{conjecture} in every sentence that states it.

\begin{center}
\begin{tabular}{ll}
\toprule
commit & \texttt{d4126d34} \\
module \texttt{sha256} & \texttt{4a6c4d680becd37acc9d02b6} \\
                       & \texttt{5e0995ff0dc26c5cd7f4edfa51468a6c04029793} \\
toolchain & \texttt{leanprover/lean4:v4.29.0-rc6} \\
Mathlib pin & \texttt{07642720480157414db592fa85b626dafb71355b} \\
elaboration & exit \texttt{0}; empty log --- zero errors, zero warnings, zero \texttt{sorry} \\
axiom oracle & 29/29 declarations on \texttt{[propext, Classical.choice, Quot.sound]} \\
\texttt{sorryAx} & none \\
\bottomrule
\end{tabular}
\end{center}

The module hash was computed independently on the author's machine and on the
verification host and agrees, so the checked object is the source printed here.
Compilation ran on a CPU (never GPU) high-memory cloud runtime, per the
project's environment rule of 1~August~2026; the desktop compiles nothing.

\textbf{One registered check remains undone.} The pre-registration also asked
that the module be imported into the project's aggregate library and that the
build job count increment by exactly one. It has not been imported, so there is
no increment to measure; measuring it honestly requires two full builds of the
aggregate in the same tree, which the shared Mathlib cache does not cover. That
check is therefore reported as \emph{not run} rather than dropped, and it is not
replaced by a comparison against a count measured elsewhere. Nothing in the
theorems above depends on it.

\begin{thebibliography}{9}
\bibitem{sylvester1852} J.~J.~Sylvester,
  \emph{A demonstration of the theorem that every homogeneous quadratic
  polynomial is reducible by real orthogonal substitutions to the form of a sum
  of positive and negative squares}, Philosophical Magazine \textbf{4} (1852),
  138--142.

\bibitem{ostrowski1959} A.~M.~Ostrowski,
  \emph{A quantitative formulation of Sylvester's law of inertia},
  Proc.\ Nat.\ Acad.\ Sci.\ U.S.A.\ \textbf{45} (1959), no.~5, 740--744.

\bibitem{ostrowski1960} A.~M.~Ostrowski,
  \emph{A quantitative formulation of Sylvester's law of inertia, II},
  Proc.\ Nat.\ Acad.\ Sci.\ U.S.A.\ \textbf{46} (1960), no.~6, 859--862.

\bibitem{hornjohnson} R.~A.~Horn and C.~R.~Johnson,
  \emph{Matrix Analysis}, 2nd ed., Cambridge University Press, 2013.
  (Congruence and inertia, \S4.5; continuity of eigenvalues, Appendix D.)

\bibitem{vandersluis1969} A.~van der Sluis,
  \emph{Condition numbers and equilibration of matrices},
  Numer.\ Math.\ \textbf{14} (1969), 14--23.

\bibitem{onsager1944} L.~Onsager,
  \emph{Crystal statistics.\ I.\ A two-dimensional model with an
  order-disorder transition}, Phys.\ Rev.\ \textbf{65} (1944), 117--149.

\bibitem{kaufman1949} B.~Kaufman,
  \emph{Crystal statistics.\ II.\ Partition function evaluated by spinor
  analysis}, Phys.\ Rev.\ \textbf{76} (1949), 1232--1243.

\bibitem{mathlib} The mathlib Community,
  \emph{The Lean mathematical library}, Proc.\ CPP 2020, ACM, 367--381.
\end{thebibliography}

\end{document}
